\documentclass[10pt,a4paper]{amsart}
\usepackage{amsmath,amssymb,amsthm,graphicx}
\usepackage{verbatim}
\newenvironment{reference}{\comment}{\endcomment}
%\newenvironment{reference}{}{}
\newenvironment{slogan}{\comment}{\endcomment}
\newenvironment{history}{\comment}{\endcomment}

% For commutative diagrams we use Xy-pic
\usepackage[all]{xy}

% We use 2cell for 2-commutative diagrams.
\xyoption{2cell}
% Theorem environments.
%
\theoremstyle{plain}
\newtheorem{theorem}[subsection]{Theorem}
\newtheorem{proposition}[subsection]{Proposition}
\newtheorem{lemma}[subsection]{Lemma}

\theoremstyle{definition}
\newtheorem{definition}[subsection]{Definition}
\newtheorem{example}[subsection]{Example}
\newtheorem{exercise}[subsection]{Exercise}
\newtheorem{situation}[subsection]{Situation}

\theoremstyle{remark}
\newtheorem{remark}[subsection]{Remark}
\newtheorem{remarks}[subsection]{Remarks}

\numberwithin{equation}{subsection}

% Macros
%
\def\lim{\mathop{\mathrm{lim}}\nolimits}
\def\colim{\mathop{\mathrm{colim}}\nolimits}
\def\Spec{\mathop{\mathrm{Spec}}}
\def\Hom{\mathop{\mathrm{Hom}}\nolimits}
\def\Ext{\mathop{\mathrm{Ext}}\nolimits}
\def\SheafHom{\mathop{\mathcal{H}\!\mathit{om}}\nolimits}
\def\SheafExt{\mathop{\mathcal{E}\!\mathit{xt}}\nolimits}
\def\Sch{\mathit{Sch}}
\def\Mor{\mathop{\mathrm{Mor}}\nolimits}
\def\Ob{\mathop{\mathrm{Ob}}\nolimits}
\def\Sh{\mathop{\mathit{Sh}}\nolimits}
\def\NL{\mathop{N\!L}\nolimits}
\def\CH{\mathop{\mathrm{CH}}\nolimits}
\def\proetale{{pro\text{-}\acute{e}tale}}
\def\etale{{\acute{e}tale}}
\def\QCoh{\mathit{QCoh}}
\def\Ker{\mathop{\mathrm{Ker}}}
\def\Im{\mathop{\mathrm{Im}}}
\def\Coker{\mathop{\mathrm{Coker}}}
\def\Coim{\mathop{\mathrm{Coim}}}

% Boxtimes
%
\DeclareMathSymbol{\boxtimes}{\mathbin}{AMSa}{"02}

%
% Macros for moduli stacks/spaces
%
\def\QCohstack{\mathcal{QC}\!\mathit{oh}}
\def\Cohstack{\mathcal{C}\!\mathit{oh}}
\def\Spacesstack{\mathcal{S}\!\mathit{paces}}
\def\Quotfunctor{\mathrm{Quot}}
\def\Hilbfunctor{\mathrm{Hilb}}
\def\Curvesstack{\mathcal{C}\!\mathit{urves}}
\def\Polarizedstack{\mathcal{P}\!\mathit{olarized}}
\def\Complexesstack{\mathcal{C}\!\mathit{omplexes}}
% \Pic is the operator that assigns to X its picard group, usage \Pic(X)
% \Picardstack_{X/B} denotes the Picard stack of X over B
% \Picardfunctor_{X/B} denotes the Picard functor of X over B
\def\Pic{\mathop{\mathrm{Pic}}\nolimits}
\def\Picardstack{\mathcal{P}\!\mathit{ic}}
\def\Picardfunctor{\mathrm{Pic}}
\def\Deformationcategory{\mathcal{D}\!\mathit{ef}}
\usepackage{cleveref}
\crefname{originallemma}{original lemma}{original lemmas}
\Crefname{originallemma}{Original lemma}{Original lemmas}
\makeatletter
\newlabel{lemma-base-change-universally-closed}{{03IS}{source}{}{sourcepointer-1}{}}
\newlabel{lemma-base-change-universally-closed@cref}{{[originallemma][0][]03IS}{[1][0][]source}{}{}{}}
\makeatother
\crefname{originallemma}{original lemma}{original lemmas}
\Crefname{originallemma}{Original lemma}{Original lemmas}
\makeatletter
\newlabel{lemma-compose-after-separated}{{03KR}{source}{}{sourcepointer-2}{}}
\newlabel{lemma-compose-after-separated@cref}{{[originallemma][0][]03KR}{[1][0][]source}{}{}{}}
\makeatother
\crefname{originallemma}{original lemma}{original lemmas}
\Crefname{originallemma}{Original lemma}{Original lemmas}
\makeatletter
\newlabel{lemma-composition-universally-closed}{{03IU}{source}{}{sourcepointer-3}{}}
\newlabel{lemma-composition-universally-closed@cref}{{[originallemma][0][]03IU}{[1][0][]source}{}{}{}}
\makeatother
\crefname{originallemma}{original lemma}{original lemmas}
\Crefname{originallemma}{Original lemma}{Original lemmas}
\makeatletter
\newlabel{lemma-permanence-finite-type}{{0462}{source}{}{sourcepointer-4}{}}
\newlabel{lemma-permanence-finite-type@cref}{{[originallemma][0][]0462}{[1][0][]source}{}{}{}}
\makeatother
\crefname{originallemma}{original lemma}{original lemmas}
\Crefname{originallemma}{Original lemma}{Original lemmas}
\makeatletter
\newlabel{lemma-quasi-compact-permanence}{{03KS}{source}{}{sourcepointer-5}{}}
\newlabel{lemma-quasi-compact-permanence@cref}{{[originallemma][0][]03KS}{[1][0][]source}{}{}{}}
\makeatother
\crefname{originallemma}{original lemma}{original lemmas}
\Crefname{originallemma}{Original lemma}{Original lemmas}
\makeatletter
\newlabel{lemma-semi-diagonal}{{03KO}{source}{}{sourcepointer-6}{}}
\newlabel{lemma-semi-diagonal@cref}{{[originallemma][0][]03KO}{[1][0][]source}{}{}{}}
\makeatother
\title{Universal closedness and properness are preserved by the factorization through a separated base}
\author{The Stacks Project Authors}
\date{stacks-project:spaces-morphisms.tex:local-snapshot-20261012; GNU FDL 1.2 or later}
\begin{document}
\maketitle
\setcounter{subsection}{1}
\noindent\textit{Editorial note.} All delivered mathematical text is the author's own, copied from the pinned Stacks Project chapter file spaces-morphisms.tex; no statement or proof step was rewritten, reordered or removed.\par\medskip
\noindent\textit{Editorial note.} The proof displays the commutative triangle as part of the statement and refers to it by X, Y and B; the factorization used in the proof is written out.\par\medskip
\noindent\textit{Editorial note.} The six references of the counted spans are recorded as non-load-bearing navigation pointers; the argument that consumes them stays inside the counted proof span.\par\medskip
\noindent\textit{Editorial note.} Preamble: the three native excerpt spans of sources/stacks-spaces-morphisms/preamble.tex (lines 8-13, 14-18 and 154-217). The chapter file itself carries no preamble text.\par\medskip
\noindent\textit{Allowed standard prerequisites (editorial).}\begin{itemize}
\item The graph (or semi-diagonal) is a closed immersion when the target is separated. Lemma lemma-semi-diagonal, used at proof lines 8193-8196. Makes the first factor of X -> Y a closed immersion, hence universally closed.
\item Universal closedness is stable under base change. Lemma lemma-base-change-universally-closed, used at proof lines 8197-8200. Makes the projection X x\_B Y -> Y universally closed.
\item A composite of universally closed morphisms is universally closed. Lemma lemma-composition-universally-closed, used at proof lines 8201-8203. Concludes part (1).
\item Permanence statements for separatedness, quasi-compactness and finite typeness. Lemmas lemma-compose-after-separated, lemma-quasi-compact-permanence and lemma-permanence-finite-type, used at proof lines 8204-8207. Upgrade the conclusion of (1) to properness in part (2).
\end{itemize}\medskip
\ifdefined\hypertarget\hypertarget{sourcepointer-1}{}\fi
\noindent\textit{Source reference.} Original lemma 03IS: focused-seeds/AGP/sources/stacks-snapshot-20261011/tags.txt, source lines 4614--4614. Navigation pointer into the Stacks Project; the standard fact it supplies is applied inside the counted proof span and the argument resting on it stays there.\par\medskip
\ifdefined\hypertarget\hypertarget{sourcepointer-2}{}\fi
\noindent\textit{Source reference.} Original lemma 03KR: focused-seeds/AGP/sources/stacks-snapshot-20261011/tags.txt, source lines 4687--4687. Navigation pointer into the Stacks Project; the standard fact it supplies is applied inside the counted proof span and the argument resting on it stays there.\par\medskip
\ifdefined\hypertarget\hypertarget{sourcepointer-3}{}\fi
\noindent\textit{Source reference.} Original lemma 03IU: focused-seeds/AGP/sources/stacks-snapshot-20261011/tags.txt, source lines 4616--4616. Navigation pointer into the Stacks Project; the standard fact it supplies is applied inside the counted proof span and the argument resting on it stays there.\par\medskip
\ifdefined\hypertarget\hypertarget{sourcepointer-4}{}\fi
\noindent\textit{Source reference.} Original lemma 0462: focused-seeds/AGP/sources/stacks-snapshot-20261011/tags.txt, source lines 5457--5457. Navigation pointer into the Stacks Project; the standard fact it supplies is applied inside the counted proof span and the argument resting on it stays there.\par\medskip
\ifdefined\hypertarget\hypertarget{sourcepointer-5}{}\fi
\noindent\textit{Source reference.} Original lemma 03KS: focused-seeds/AGP/sources/stacks-snapshot-20261011/tags.txt, source lines 4688--4688. Navigation pointer into the Stacks Project; the standard fact it supplies is applied inside the counted proof span and the argument resting on it stays there.\par\medskip
\ifdefined\hypertarget\hypertarget{sourcepointer-6}{}\fi
\noindent\textit{Source reference.} Original lemma 03KO: focused-seeds/AGP/sources/stacks-snapshot-20261011/tags.txt, source lines 4684--4684. Navigation pointer into the Stacks Project; the standard fact it supplies is applied inside the counted proof span and the argument resting on it stays there.\par\medskip
\begin{lemma}
\label{lemma-universally-closed-permanence}
Let $S$ be a scheme.
Consider a commutative diagram of algebraic spaces
$$
\xymatrix{
X \ar[rr] \ar[rd] & &
Y \ar[ld] \\
& B &
}
$$
over $S$.
\begin{enumerate}
\item If $X \to B$ is universally closed and $Y \to B$ is
separated, then the morphism $X \to Y$ is universally closed.
In particular, the image of $|X|$ in $|Y|$ is closed.
\item If $X \to B$ is proper and $Y \to B$ is separated, then
the morphism $X \to Y$ is proper.
\end{enumerate}
\end{lemma}

\begin{proof}
Assume $X \to B$ is universally closed and $Y \to B$ is separated.
We factor the morphism as $X \to X \times_B Y \to Y$.
The first morphism is a closed immersion, see
Lemma \ref{lemma-semi-diagonal}
hence universally closed.
The projection $X \times_B Y \to Y$ is the base change
of a universally closed morphism and hence
universally closed, see
Lemma \ref{lemma-base-change-universally-closed}.
Thus $X \to Y$ is universally closed as the composition
of universally closed morphisms, see
Lemma \ref{lemma-composition-universally-closed}.
This proves (1). To deduce (2) combine (1) with
Lemmas \ref{lemma-compose-after-separated},
\ref{lemma-quasi-compact-permanence}, and
\ref{lemma-permanence-finite-type}.
\end{proof}

\end{document}
